\scanpage{41}
\begin{center}- 18 -\end{center}
\hfill n$^\circ$ 276

un module filtré $(E,E_n)$ sur un anneau filtré $(A,A_n)$, et soit $\hat E$ le complété du groupe abélien $E$ muni de la topologie associée à la filtration donnée $\hat E$ est donc un groupe topologique abélien, d'ailleurs métrisable puisque la topologie de $E$ est définie par un écart. Comme les sous-groupes $E_n$ forment un système fondamental de voisinages de 0 dans $E$, \uline{les sous-groupes} $\hat E_n$ \uline{forment de même un système fondamental de voisinages de} 0 \uline{dans le groupe} $\hat E$ et si $F$ est un sous-groupe de $E$, muni de la topologie induite, le complété de $E/F$ n'est autre que $\hat E/\hat F$. Prenant $F=E_n$ et observant que, $E_n$ étant ouvert, le quotient $E/E_n$ est discret et donc complet, on voit que \uline{les homomorphismes canoniques}
\[
E/E_n\longrightarrow\hat E/\hat E_n
\]
\uline{sont bijectifs}. On notera d'autre part que, $E_n$ étant fermé, on a, lorsque $E$ est séparé (i.e. se plonge biunivoquement dans $\hat E$), la relation.
\[
E_n=E\cap\hat E_n.
\]
Considérons plus particulièrement l'anneau filtré $(A,A_n)$ lui-même ; Alors $\hat A$ est un \uline{anneau topologique} (\uline{Top.Gén.} Chap.III, \S\ 5, n$^\circ$4) et les relations $A_pA_q\subset A_{p+q}$ conduisent aux relations $\hat A_p\hat A_q\subset\hat A_{p+q}$ ; les sous-groupes $\hat A_p$ définissent donc une filtration de l'anneau $\hat A$, et comme ils forment un système fondamental de 0 on voit que la topologie de $A$ est précisément définie par cette filtration. Autrement dit, \uline{le complété de l'anneau filtré} $(A,A_n)$ \uline{est l'anneau filtré} $(\hat A,\hat A_n)$. Comme de plus $\hat A/\hat A_n$ s'identifie canoniquement à $A/A_n$ il est clair que \uline{les anneaux gradués associés} $G(A)$ \uline{et} $G(\hat A)$ \uline{sont canoniquement isomorphes}.

Revenant à un module filtré $(E,E_n)$ sur l'anneau filtré $(A,A_n)$, la loi de composition $A\times E\longrightarrow E$ se prolonge en une loi de composition $\hat A\times\hat E\longrightarrow\hat E$ on vérifie immédiatement qu'elle fait de $\hat E$ un $\hat A$-module ; comme de plus les
\mspage{40}
